    \documentclass[10pt]{article}

\def\Type{LessonDJ}
\def\TypeNum{Workshop Week 8} %Lesson #
\def\TypeTitle{More 1st and 2nd  Derivatives+Optimization } %Lesson Title
\def\Semester{Fall 2021} %Not Used 
\def\Course{MATH 1190}%Course number and name
\def\Targets{  \target{A2}, \target{A3} }

\usepackage{preambleDJ} 

%%%%%%%%%%%% BEGIN DOCUMENT %%%%%%%%%%%%%%%
%%%%%%%%%%%% BEGIN DOCUMENT %%%%%%%%%%%%%%%
%%%%%%%%%%%% BEGIN DOCUMENT %%%%%%%%%%%%%%%
%%%%%%%%%%%% BEGIN DOCUMENT %%%%%%%%%%%%%%%
%%%%%%%%%%%% BEGIN DOCUMENT %%%%%%%%%%%%%%%

\begin{document}
\pagestyle{otherpages}
\thispagestyle{firstpage}

\vs{.6}
\
%\begin{center}
%    \includegraphics[width=0.8\textwidth]{IMG_2891.jpeg}
%\end{center}
\section{Local and Absolute Max/Mins and Concavity}
\begin{enumerate}

\item \ltrm{\target{A2}}  Given $\ds f(x) =  \frac{x^2-2x+2}{x-1}. $
\enum{
\item Find $f^{\prime}(x)$ and determine all critical numbers (there are a total of $3$).
\item Construct a sign chart for $f'$ and determine which intervals the function is increasing and decreasing
\item For each critical number, decide if it is a local minimum, local maximum, or neither.
\item $\ds f^{\prime\prime}(x)=\frac{2}{(x-1)^3}$. Create a sign chart for $f''$.  Use this to find the intervals where $f$ is concave up and concave down. Are there any inflection points? Why or why not?
\item For each  interval, find the absolute maximum and minimum (if it exists).  
\tasksR{2}{
\task  $[2,5]$
\task  $(-\infty,1)$ 
}
\item Sketch the graph of $f(x)$. Note that there is a vertical asymptote at $x=1$.
}




\newpage
\section{More Local and Absolute Max/Mins and Concavity}

\item \ltrm{\target{A2}} Given $f(x)= \ln(x-1)-x$.
\begin{enumerate}
    \item What is its domain? Recall that $\ln(x)$ is only defined for $x > 0$. 
    \item Find the critical points of $f(x)= \ln(x-1)-x$ (there are two). Do each step in detail. 
    \item Create a sign chart for $f'$. On what intervals is this function increasing, on which is it decreasing? 
    \item Determine if the critical points locate local maxima, local minima, or neither.
    \item When is $f(x)$ concave up? When is $f(x)$ concave down? Any inflection points?
    \item For each  interval, find the absolute maximum and minimum (if it exists).  
    \tasksR{2}{
    \task $(1,5)$
    \task $[2,e+1]$
    }
    \item Sketch a graph of $f(x)= \ln(x-1)-x$.

    \end{enumerate}

\newpage
\section{Optimization}
\item \ltrm{\target{A2}} Given $\ds f(x) =  15x-x^2. $
\enum{
\item Find $f^{\prime}(x)$ and determine all critical numbers.
\item Construct a sign chart for $f'$ and determine which intervals the function is increasing and decreasing.
\item For each critical number, decide if it is a local minimum, local maximum, or neither.
\item Find the Absolute Max/Min for the interval $(0, \infty)$
}

\mpageT{.5}{.5}{
\item \ltrm{\target{A3}}  Suppose Bob wants to build a fence to enclose some sheep with a rectangular sheep pen. Fortunately, he lives next to a river, so he only needs to build three sides of the enclosure. Bob has $30$m of fencing. What are the dimensions that will enclose the greatest amount of sheep?
}{\graphicsC{sheepPen}{3}
}
{\bf Modeling Phase}
\tasksA{2}{
\task The variables $x$ and $y$ describe the width and length of the sheep pen. variables that describe the width and length of the sheep pen.  What quantity are we going to maximize? Let's call this quantity $S$.  Use both of your  variables to write down an equation for this quantity.\vs{1}
\task Bob's constraint is that he only has $30$m of fencing. Use your variables to write down an equation that represents this constraint.
\task Both the quantity you are going to maximize, $S$, and your constraint contain both of your variables. Use your constraint to solve for one variable and substitute  into $S$.
}
\vs{2}
{\bf Calculus Phase} Think about the process you have been using to optimize other quantities. Use that process here to find the dimension that maximize the area of the sheep pen? {\em Hint:} Consider using work you have done on an earlier problem.

%
%\newpage
%\item  Given $\ds f(x) =  \frac{3-x}{8}+\frac{\sqrt{x^2+1}}{5}. $
%\enum{
%\item Find $f^{\prime}(x)$ and use technology (e.g. Desmos) to determine all critical numbers.
%\item Construct a sign chart for $f'$ and determine which intervals the function is increasing and decreasing.
%\item For each critical number, decide if it is a local minimum, local maximum, or neither.
%
%\item Suppose John is standing next to a river that is 1 mile wide. His house is on the opposite of the river 3 miles down. John can swim at 5 miles per hour but he can run at 8 miles per hour. What is the shortest amount of time it can take John to reach his house? You will find this by finding the best choice for $d$ where John stops running and starts swimming.\\
%\mpageT{.4}{.6}[\hs{.1}]{
%{\bf Modeling Phase} We first need to develop a model of the total time: $$T_{total}=T_{run}+T_{swim}$$
%The total horizontal distance to the house is $3$ miles.  Once John starts to swim, the remaining horizontal distance is $d$ miles (see Figure 1). 
%}{\graphicsC{runSwim}{4}{Figure 1}
%}
%\tasksA{2}{
%\task What is the distance John ran? Your answer should contain $d$.
%\task How long does it take for John to run this distance? This is $T_{run}$.  Your answer should contain $d$.\vs{.5}
%\task What is the distance John swam?  {\em Hint}: You will want to utilize Pythagorean's Theorem to find this distance. Your answer should contain $d$.
%\task How long does it take for John to swim this distance? This is $T_{swim}$.  Your answer should contain $d$. \vs{.75}
%\task* What is $T_{total}$? Your answer will contain $d$\vs{.25}
%}
%{\bf Calculus Phase} Think about the process you have been using to optimize quantities. Use that process to find the minimum time it takes for John to get home.
%}
%
%    \vfill













\end{enumerate}

\end{document}